\label{cap:introducao_modelo}

O presente Trabalho de Conclusão de Curso foi desenvolvido em regime de cooperação técnico-científica, em coautoria com o acadêmico Victor de Souza Vilela da Silva, do Curso de Engenharia de Computação do Centro Federal de Educação Tecnológica de Minas Gerais, Unidade Leopoldina. Esta pesquisa constitui-se como parte complementar e interdependente do estudo intitulado \textit{Interpretador Dinâmico para Ensino de Programação: Implementação de análise léxica e sintática, geração e interpretação de código intermediário em ambiente web}, de autoria do referido coautor e submetido ao mesmo curso desta Instituição.

Embora ambos os documentos partilhem a mesma fundamentação macro e o mesmo objeto de estudo, a saber, uma plataforma educacional web composta por um interpretador dinâmico acoplado a uma IDE com gramática configurável pelo usuário, as abordagens diferem em escopo e ênfase metodológica. Enquanto o documento complementar se concentra no núcleo de tradução, examinando as fases de análise léxica, análise sintática, geração de código intermediário e execução pelo interpretador, este volume volta-se ao Ambiente de Desenvolvimento Integrado: sua arquitetura front-end, a concepção e implementação do assistente de personalização da linguagem, a experiência de uso da plataforma e os procedimentos de coleta de retroalimentação junto aos estudantes. Para uma compreensão holística e integral do sistema proposto, recomenda-se a leitura cruzada e integrada de ambos os documentos.

As reflexões que abrem este capítulo, bem como as considerações sobre as dificuldades enfrentadas por estudantes em seu primeiro contato com a programação, apoiam-se principalmente nos estudos de \citeonline{stefik2013} e \citeonline{ferreira2025}.

A familiaridade mínima com a programação deixou de ser uma exigência restrita aos cursos de Ciência da Computação e da Engenharia de Computação. Áreas tão distintas quanto economia, biologia, ciências sociais e jornalismo passaram a incorporar a manipulação programática de dados como ferramenta de trabalho corrente, e a tendência observável é de aprofundamento dessa transversalidade. Em contraste com esse cenário, as disciplinas de iniciação à programação seguem registrando taxas elevadas de reprovação e de evasão, fenômeno cujo enfrentamento exige compreender, com algum cuidado, em que ponto da curva de aprendizagem o estudante costuma encontrar maior resistência.

Observa-se, recorrentemente, que a barreira não está no raciocínio algorítmico em si. Estudantes que demonstram capacidade de descrever, em linguagem natural, a sequência de passos necessária para resolver um problema esbarram, no momento de transpor esse raciocínio para o computador, na rigidez das linguagens de programação tradicionais. Pontos e vírgulas omitidos, chaves desbalanceadas, palavras reservadas escritas em idioma estrangeiro e mensagens de erro hermeticamente formuladas formam um conjunto de pequenas exigências que, embora secundárias do ponto de vista lógico, consomem parcela considerável do tempo e da energia cognitiva do iniciante. A frustração resultante desse atrito tende a se acumular nas primeiras semanas de curso, justamente o período em que a percepção de progresso é mais decisiva para a continuidade do estudo.

A pergunta que orienta esta pesquisa nasce dessa observação: \textbf{é possível conceber uma plataforma de programação que permita ao estudante, ele próprio, redefinir o vocabulário e as regras estruturais da linguagem com que aprende, de modo a deslocar o foco da forma para o conteúdo do raciocínio computacional?}

A hipótese defendida ao longo do trabalho responde afirmativamente a essa indagação. Conjectura-se que, ao colocar nas mãos do discente as decisões sobre quais palavras servirão como comandos, sobre como blocos serão delimitados, sobre como os operadores serão escritos e sobre o caráter obrigatório ou facultativo do terminador de instrução, reduz-se de maneira apreciável a carga cognitiva exigida no aprendizado simultâneo da lógica e da forma. O exercício de adaptar a linguagem ao próprio repertório transforma a sintaxe, antes vivida como obstáculo, em objeto deliberado de escolha pedagógica.

\section{Justificativa}

A motivação para empreender este projeto repousa sobre uma constatação proveniente da análise das ferramentas educacionais hoje disponíveis: ainda que diversas plataformas se proponham a facilitar o ensino de programação, a esmagadora maioria opera sob uma gramática fixa, definida \textit{a priori} por seus mantenedores. Linguagens didáticas em português, juízes online voltados à prática contínua e ambientes de programação visual têm méritos próprios; nenhum deles, contudo, oferece ao usuário final a faculdade de reescrever as regras estruturais que regem a forma como o pensamento é convertido em código.

A construção de uma ferramenta que inverta essa relação, adaptando a linguagem ao estudante em vez de exigir que o estudante se adapte à linguagem, é assumida aqui como uma contribuição inédita no campo da informática educacional. Diferentemente de iniciativas voltadas estritamente ao núcleo de tradução ou estritamente à interface, este trabalho aposta no acoplamento dessas duas frentes: um interpretador capaz de aceitar gramáticas configuradas pelo usuário e uma IDE web cuja experiência de uso oferece, ao mesmo tempo, conforto operacional próximo ao de ferramentas profissionais e mediações pedagógicas adequadas ao público iniciante.

\section{Objetivos}

A seguir são apresentados o objetivo geral que orienta a investigação e os objetivos específicos que viabilizam o seu alcance.

\subsection{Objetivo Geral}

Construir uma plataforma educacional formada por uma IDE web e por um interpretador dinâmico que aceite, em tempo de execução, regras gramaticais definidas pelo próprio usuário, a fim de oferecer suporte ao ensino e à aprendizagem de algoritmos por estudantes em estágio introdutório.

\subsection{Objetivos Específicos}

Para que o objetivo geral seja efetivamente alcançado, as seguintes metas operacionais foram fixadas:

\begin{itemize}
    \item Estudar, sob perspectiva crítica, os fundamentos de construção de compiladores e as ferramentas educacionais disponíveis na literatura;
    \item Participar do projeto e da integração do interpretador, contemplando a execução das análises léxica e sintática e do código intermediário a partir da IDE;
    \item Conceber a IDE web, dotando-a de recursos para personalização da linguagem, execução de código no próprio ambiente e apresentação de \textit{feedback} ao usuário;
    \item Validar a plataforma por meio de estudos de caso em ambiente educacional, observando aspectos de usabilidade, de viabilidade técnica e de potencial pedagógico.
\end{itemize}

\section{Estrutura do Trabalho}

O restante deste documento desdobra-se da seguinte maneira. O Capítulo~\ref{cap:referencial} percorre o referencial teórico que sustenta o projeto, articulando os fundamentos de construção de compiladores, as práticas associadas ao desenvolvimento de aplicações web e as iniciativas educacionais consideradas correlatas. O Capítulo~\ref{cap:proposta}, por sua vez, expõe a proposta do trabalho e o cronograma seguido em suas duas fases. Por fim, o Capítulo~\ref{cap:modelagem} descreve a modelagem efetivamente realizada e a implementação consolidada até o encerramento desta etapa, contemplando a arquitetura da plataforma, a personalização dinâmica da linguagem, a integração com o interpretador e a estratégia de testes adotada.
